\documentclass[conference]{IEEEtran}
\usepackage{booktabs}
\usepackage{amsmath}
\usepackage{amssymb}
\usepackage{graphicx}
\usepackage{xurl}
\usepackage{balance}
\usepackage{etoolbox}
\usepackage[hidelinks]{hyperref}

\makeatletter
\patchcmd{\abstract}{---}{: }{}{}
\patchcmd{\IEEEkeywords}{---}{: }{}{}
\makeatother

\title{Silent Adoption or Non-Commitment? Knowledge-Base Poisoning and Record-Scope Costs in Utility-Operations RAG}

\author{
\IEEEauthorblockN{Nassim Tinkicht, Hussain Al-Aqrabi}
\IEEEauthorblockA{Computer Information Science, Higher Colleges of Technology, United Arab Emirates\\
Email: ntinkicht@hct.ac.ae; halaqrabi@hct.ac.ae}
}

\begin{document}
\bstctlcite{GridPoisonBSTControl}
\maketitle

\begin{abstract}
Retrieval-augmented generation (RAG) can ground utility lookups in documents, but indexed repositories can admit low-trust content.
We introduce GridPoison-100: 35 public-standard, 45 fictional asset-register, and 20 legitimate scope-conflict items, and test targeted numeric poisoning on three LLMs.
Both pre-specified primary tests were null: in a fixed context containing three clean same-item supports plus one poison, no model adopted the attacker value, so the public-versus-asset comparison was uninformative; and a keyed system-of-record (SoR) upper-bound check at 70\% nominal coverage was not significantly better than a cautious prompt or simplified isolate-then-aggregate voting (all Holm-adjusted $p=1.0$; significance was unattainable with the observed sparse discordance).
Descriptively, one dense-retrieved poison at $k=5$ reached 49/70 public-query trials per model and reduced correct answers from 75.7--78.6\% to 8.6--12.9\%; adoption was 2/49 and failures were mostly non-commitment.
Because each poison is a minimal pair of one clean support, this low adoption is not general evidence that clean context protects: locked forced-context follow-ups produced 0/80 adoption with the paired twin, entirely through non-commitment, 15--17/80 with a different clean support, and 80/80 with the poison alone.
A locked commit-required prompt raised adoption to 24/49 for Mistral Small 3.2 and 8/49 for GPT-5.6 Luna. Mistral met the pre-stated material-effect band; Luna's point estimate fell in the partial band, but its interval crossed all three bands and 12/49 outputs were empty.
Full keyed record coverage removed observed out-of-range commitments by construction; deliberately generic records for 13/20 scope-conflict items false-flagged 64.1--65.7\% of otherwise-correct answers.
These results separate silent integrity loss from abstention and illustrate the need for scope-aware authoritative records.
\end{abstract}

\begin{IEEEkeywords}
retrieval-augmented generation, knowledge-base poisoning, power systems, utility operations, large language models, data integrity
\end{IEEEkeywords}

\section{Introduction}
Large language models are increasingly proposed as interfaces to technical documentation, including power-system applications~\cite{ruan2024power,lee2025smartgridrag}. A plausible near-term utility use is constrained lookup: workers retrieve approach distances, interconnection boundaries, or approved asset limits. RAG makes such answers document-grounded but also creates a write-path dependency: low-trust contributions from contractors, shared drives, wikis, or external bulletins can enter the index.

Knowledge-base poisoning is already established~\cite{zou2025poisonedrag,pereira2026influence}; our question is how a conflict-aware utility assistant separates silent corruption from non-commitment, and how structured-record validation trades filtering benefit against false flags on legitimate scope conflicts.

\textbf{RQ1} asks how poisoning decomposes into retrieval exposure and conditional adoption. \textbf{RQ2} compares public-standard and asset-specific adoption at fixed exposure. \textbf{RQ3} compares keyed SoR validation with two generic defenses and measures false flags when legitimate documents differ by jurisdiction or interconnection. Contributions are: (1) GridPoison-100 and a utility write-path threat model; (2) a locked commit-required intervention that distinguishes prompt-mediated non-commitment from silent adoption; (3) a pre-stated composition test contrasting paired clean twins with non-paired supports; and (4) record-validation accounting that exposes the tradeoff between filtering and benign-conflict false flags under incomplete or generic record scope.

\section{Related Work}
PoisonedRAG established targeted answer steering through retrieved attacker text~\cite{zou2025poisonedrag}; later work studied retrieval, database, and generator factors affecting poisoning~\cite{pereira2026influence}. Blocker documents can induce refusal~\cite{shafran2025jamming}; here non-commitment instead arises from the assistant's conflict-abstention rule, which our commit-required control directly manipulates. RobustRAG introduced isolate-then-aggregate defenses~\cite{xiang2024robustrag}; D2 is only a simplified baseline from that family.

ClashEval and controlled conflict studies show that retrieved evidence can compete with model priors and that contradictory evidence can change answer or abstention behavior~\cite{wu2024clasheval,xie2024chameleon,han2026evidenceconflicts}. Evidence phrasing, recency, and metadata also affect preference~\cite{wan2024evidence,chiang2024metadata}. Power-system LLM security risks and smart-grid threat models have already been studied~\cite{ruan2024power,li2024smartgridrisk}, while smart-grid RAG prototypes have targeted secure offline document retrieval~\cite{lee2025smartgridrag}. Our domain contribution is therefore not RAG deployment itself. Utility specificity comes from versioned protection curves and state DER rules, asset-register authority, contractor/bulletin write paths, and jurisdiction/interconnection scope; we use these to measure silent corruption, prompt-dependent non-commitment, and record-scope false flags.

\section{Threat Model}
The attacker is an authorized but low-trust writer, such as a contractor or bulletin sender whose content is indexed. They may add one to three documents but cannot edit existing documents, change the structured record, EMS/SCADA, protection devices, model weights, or the user's query. They know the target asset/topic and existing documents, but not the exact query wording, embedding model, or realized retrieval.

The goal is targeted numeric answer-integrity corruption: a benign lookup commits to an attacker-chosen value. The primary poison is factual text rather than an instruction; one stress paraphrase inherits directive wording from a clean template while changing only the value. Confidentiality, weight poisoning, multi-turn social engineering, and control-loop attacks are outside scope. Denial of service is not the goal, but loss of correct answers is reported because it dominates observed failures.

\section{GridPoison-100}
GridPoison-100 has 100 items. Tier~1 contains 35 public facts from Massachusetts DER settings/return-to-service material, OSHA 29 CFR 1910.269 Tables R-6/R-7, and NERC PRC-024-3 boundaries~\cite{osha1910269,nercprc0243,matsrgsettings,matsrgservice}. OSHA items encode voltage range, exposure type, and elevation at or below 900~m. FERC approved PRC-024-4 and PRC-029-1 in 2025, and NERC lists Oct.~1, 2026 as their U.S. effective date. The benchmark therefore freezes PRC-024-3 as the source version used for the study and treats standard version and jurisdiction as part of scope~\cite{fercorder909,nerceffective2026}. Each public item stores source, locator, units, applicability, and an allowable interval.

Tier~2 contains 45 synthetic asset-specific values from a declared fictional register covering transformer, breaker, feeder, BESS, and cable limits. The register is authoritative by construction and does not describe a real network.

Tier~3 contains 20 unpoisoned legitimate conflicts where documents differ by jurisdiction, interconnection, or table scope while the query names the applicable context. These items measure false flags on correct context-specific answers.

The clean corpus has 400 short template-generated documents: three supports per poisonable item, 40 legitimate-conflict documents, and 120 on-topic distractors. Identifiers are explicit, so the corpus is cleaner than a real repository. Each item has a formal query and a field-style paraphrase. Each poison is a minimal pair with one support and changes only the critical value. The main variant moves outside the benchmark interval; wrong-but-in-range and dated-bulletin variants test factual error and authority framing separately. Items were encoded by the authors from the cited sources; the planned external engineer review was not completed, so no engineer-validated, certification, or operational-safety claim is made.

\section{Experimental Design}
\subsection{Reference conditions and models}
B0 is closed-book generation, B1 clean RAG, and B2 poisoned RAG with the same base prompt. The base prompt requires \texttt{VALUE: UNKNOWN} when supplied evidence is genuinely inconsistent and the query does not resolve the conflict; D1 additionally warns that retrieved sources may be mistaken or falsified and prioritizes matching jurisdiction, interconnection, asset identifier, and SoR. The locked commit-required follow-up instead requires one best numeric value and forbids \texttt{UNKNOWN}. Retrieval uses sentence-transformers/all-MiniLM-L6-v2 cosine similarity with one document per chunk. Models are Qwen3-8B (\texttt{qwen/qwen3-8b}), Mistral Small 3.2 24B (\texttt{mistralai/mistral-small-3.2-24b-instruct}), and GPT-5.6 Luna (\texttt{openai/gpt-5.6-luna}) through OpenRouter at temperature zero and \texttt{max\_tokens}=256. Provider routing was not pinned; original B2 requests were served by DeepInfra/Venice/Parasail for Mistral, Alibaba for Qwen, and OpenAI for Luna, while the completed commit follow-up used the same Mistral provider set and OpenAI for Luna. Resolved provider/model metadata, response ID, token usage, and timestamp were logged. For Luna, hidden reasoning tokens count against the 256-token response budget and some trials therefore exhausted the budget without visible answer text.

\subsection{Stage A and Stage B}
RQ1 uses natural dense retrieval with one poison and a three-poison stress condition comprising the original minimal-pair poison plus two paraphrases; Poison@$k$ is reported separately from generation. Pre-generation calibration found strongly imbalanced one-poison exposure across Tiers~1 and~2, so natural retrieval cannot support a direct cross-tier adoption comparison. RQ2 therefore uses a fixed five-passage context: all three clean same-item supports, including the poison's paired twin, one deterministic on-topic distractor that states it does not apply to the target, and one poison, counterbalanced first/last. This evidence-rich design holds exposure and context length constant and is not a natural attack-strength estimate.

\subsection{Defenses}
D1 is the cautious prompt. D2 isolates each retrieved passage, extracts one value with that prompt, ignores non-numeric outputs, and returns the plurality numeric value; no numeric votes or a tie yields \texttt{UNKNOWN}. It is a simplified isolate-then-aggregate baseline~\cite{xiang2024robustrag}. D3 deterministically post-processes the same B2 stress outputs: if the applicable item record is covered and a committed value differs, the answer is replaced by a verification flag. The applicable record key is given, so linking errors are not modeled. Records equal the benchmark target for 78/80 poisonable items; two Massachusetts undervoltage items instead store IEEE~1547-2018 Category~II defaults~\cite{ieee1547,matsrgsettings}. For 13/20 Tier~3 items, the record intentionally stores a generic alternative. D3 therefore measures an upper-bound answer-key-like check, not a new defense algorithm. SHA-256 coverage at nominal 40/70/100\% realizes 30/57/80 of 80 poisonable items and 6/14/20 Tier~3 items.

RQ3 uses the three-poison $k=5$ stress condition to avoid a retrieval floor. Secondary one-factor deviations test a dated bulletin, wrong-but-in-range poison, and $k=2$; no full factorial sweep is performed.

\section{Metrics and Statistical Analysis}
Submission-facing outcomes are correct (\textbf{C}), attacker adoption (\textbf{A}), another committed value (\textbf{W}), and non-commitment (\textbf{N}). The planned manual split of non-numeric outputs into conflict flags versus refusals was not performed, so all non-numeric outputs are pooled as N. Numeric C/A/W grading uses the required \texttt{VALUE:} line with zero numeric tolerance. Benchmark queries, targets, and poisons use the stored item unit; the parser records trailing unit text but performs no unit conversion.

Let $E_t=1$ when at least one poison is retrieved in the top-$k$, $C_t=1$ when trial $t$ commits a numeric value, $T$ be the trial set, and $i(t)$ map a trial to its benchmark item. We report Poison@$k$, conditional adoption rate (CAR), correct-response rate (CRR), non-commitment rate (NCR), and benchmark-violation rate (BVR):
\begin{equation}
\mathrm{CAR}=P(A\mid E_t=1),
\end{equation}
\begin{equation}
\mathrm{BVR}=\frac{1}{|T|}\sum_t \mathbb{1}[C_t=1 \land \hat v_t \notin [L_{i(t)},U_{i(t)}]].
\end{equation}
For Tier~3, benign-conflict false-flag rate (BFFR) is the fraction of clean-RAG-correct trials that the keyed record check flags. BVR counts only committed numeric values, so it is reported with CRR; it is a benchmark consequence metric, not a physical grid-safety measure.

The item is the inferential unit; two query phrasings are repeated observations. Defense differences use paired 95\% item-cluster bootstrap intervals (2,000 resamples); natural-retrieval CAR intervals also cluster by item. For zero-event Stage-B cells we report exact item-level Clopper--Pearson upper bounds because the bootstrap interval is degenerate; this post-lock reporting change is documented in the methods lock. The pre-specified RQ3 family compares D3@70 with D1/D2 using item-level exact McNemar tests with Holm correction over six comparisons. The implemented item collapse counts a violation only when both phrasings violate; because the lock did not specify this repeated-phrasing collapse, an either-phrasing sensitivity is also reported. B2 comparisons are descriptive. Tier~3 BFFR is conditioned on clean-RAG correctness; this is post-hoc, and unconditional full-coverage BFFR was 62.5\% for every model. Code, data, prompts, run metadata, and the methods lock are public at \url{https://github.com/NTinkicht/Research-GridPoisonRAG}; pre-generation alignment is commit \texttt{3de75c6}, and the follow-up lock is \texttt{c598f44b}.

\section{Results}
Both pre-specified primary endpoints were null: Stage-B attacker-value adoption was 0\% in both tiers (RQ2), and no D3@70 comparison with D1 or D2 was significant (RQ3).
The exposed-trial breakdown of natural retrieval below was added after these outcomes were known and is descriptive; $k=2$ was pre-specified as a secondary one-factor deviation.
% AUTO-GENERATED results will replace this placeholder.
\subsection{Results pending}
The final result tables are generated automatically after the three-model experiment matrix completes.

\subsection{Prospective Follow-up and Exploratory Additional Model}
Four follow-up checks were locked after the original results. Qwen's commit/BM25 stages completed but were not retained when Stage C failed before artifact staging; recovery also failed, so no Qwen follow-up is reported. Mistral and Luna are the completed locked follow-ups. Gemini 2.5 Flash-Lite is exploratory with a matched baseline; no locked verdict depends on it. Stage C is falsified if any locked model crosses a threshold, so its verdict is independent of Qwen.

\emph{Design.} The frozen dense retriever was audited at top-10. A commit-required prompt reused one-poison $k=5$; CAR uses the same 49 exposed Tier~1 trials, item-cluster intervals, and item-level paired inference. BM25 used lower-cased whitespace tokens and corpus-order ties at $k=2,5$. Stage~C used one formal query per item: C0, two clean supports; C1, poison+paired twin; C2, poison+different clean support; C3, poison+distractor; C4, poison alone. The strong composition hypothesis was falsified if any locked model reached $\geq20\%$ adoption in C1/C2 or $\leq50\%$ in C4.

\emph{Results.} The audit reproduced all 640 original retrieval contexts (1920 generation rows across three models) with no poison-rank mismatch (one-poison $k=5$, three-poison $k=5$, bulletin $k=5$, and one-poison $k=2$). Across both tiers at $k=2$, 9 exposed trials had no clean same-item passage and all 9 adopted; 20 exposed trials had at least one clean same-item passage and none adopted. Because the adoption outcomes were known before this audit, this is a consistency check rather than an independent test. In the one-poison $k=5$ Tier~1 contexts, the paired clean twin was present in 47/49 exposed trials; at $k=2$ it was present in 17/26. The $k=2$ exposed set spans 19 items and the nine adopting trials span 9 items. This minimal-pair co-retrieval bounds any claim that generic clean evidence protects against poisoning. Requiring commitment raised adoption from 2/49 to 24/49 for Mistral (49.0\%; item-cluster 95\% CI 33.3--64.4; item-level exact $p=3.05\times 10^{-5}$, $n=31$ items) and from 2/49 to 8/49 for Luna (16.3\%; 7.5--26.5; item-level exact $p=0.0312$, $n=31$ items). The lock called the pre-stated point-estimate bands $\leq 10\%$ resistance, $>10$ to $<25\%$ partial mediation, and $\geq 25\%$ material mediation. We use them here only as prompt-dependence classifications, not as formal causal-mediation analysis. Mistral meets the material band. Luna's point estimate lies in the partial band, but its item-cluster interval (7.5--26.5\%) crosses all three bands. Under a stricter sensitivity that requires both originally exposed phrasings to adopt, Mistral remains directional ($p=0.0312$) while Luna has no positive discordant item ($p=1.0000$). Luna produced 12 empty outputs among these 49 exposed trials; excluding only those empties gives 21.6\% adoption. Pooled clean-control CRR was 83.8\% (Mistral), 86.9\% (Luna), and 84.4\% (Gemini, exploratory). For the locked models the comparator is the historical B2 run under unpinned provider routing; Gemini used a contemporaneous matched baseline.
\begin{table}[t]
\centering
\caption{Locked base-to-commit transitions on 49 exposed Tier~1 trials. N is non-commitment; Other contains transitions not shown separately.}
\label{tab:commit-transitions}
\footnotesize
\begin{tabular}{lrrrrr}
\toprule
Model & A$\to$A & N$\to$A & N$\to$C & N$\to$N & Other \\
\midrule
Mistral & 2 & 22 & 23 & 0 & 2 \\
Luna & 2 & 6 & 26 & 12 & 3 \\
\bottomrule
\end{tabular}
\end{table}

Empty-output audit: Mistral and Qwen had none in B2\_P1/Stage B; Luna had 1/160 in B2\_P1, 4/80 poison-first, none poison-last, and 2/80 in each of C2/C3 (none in C0/C1/C4).
B2 providers were DeepInfra/Parasail/Venice (Mistral), Alibaba (Qwen), and OpenAI (Luna); commit providers were DeepInfra/Parasail/Venice and OpenAI. Closed-book Tier~1 correct counts were 0/70, 2/70, and 0/70; Luna had 41/70 empties.
Stage~C falsified the pre-stated strong composition hypothesis: Mistral adopted in 17/80 C2 trials (21.2\%; exact 95\% CI 12.9--31.8), one trial above the threshold, while Luna (15/80) met every threshold. Both locked models adopted in 0/80 C1 and 80/80 C4 trials. In C1 the paired twin protected integrity by producing non-commitment rather than correct answers: both locked models were 0/80 correct. The no-poison C0 control was 0/80 adoption and 80/80 correct for every model. C2/C3/C4 correct counts were 0/80 for both locked models; Table~\ref{tab:upgrade-controls} reports adoption counts.
\begin{table}[t]
\centering
\caption{Follow-up controls (counts). Commit: attacker adoptions among 49 exposed Tier~1 trials under the commit-required prompt (base prompt: 2, 2, and 3). BM25: adoptions among trials exposed at $k=5$ under the locked corpus-order tie rule. C1--C4: adoptions out of 80 forced-context trials. Gemini 2.5 Flash-Lite is exploratory.}
\label{tab:upgrade-controls}
\footnotesize
\setlength{\tabcolsep}{3.5pt}
\begin{tabular}{lrrrrrr}
\toprule
Model & Commit & BM25 & C1 & C2 & C3 & C4 \\
\midrule
Mistral Small 3.2 & 24/49 & 2/38 & 0 & 17 & 53 & 80 \\
GPT-5.6 Luna & 8/49 & 0/38 & 0 & 15 & 73 & 80 \\
Gemini (expl.) & 21/49 & 12/38 & 1 & 63 & 78 & 80 \\
\bottomrule
\end{tabular}
\end{table}

Under BM25, the poison and its paired clean twin had exactly equal scores in 160/160 queries. Exposure was therefore highly tie-order dependent: under the locked clean-first corpus order, the poison appeared in 2/160 queries at $k=2$ and 38/160 at $k=5$; placing the poison first among equal scores changes those counts to 112/160 and 124/160, respectively. BM25 is therefore a tie-order sensitivity, not evidence of lexical robustness.
\emph{Post-hoc benchmark consequence mapping.} A post-hoc benchmark-consequence mapping decomposed each OSHA MAD into its electrical component and ergonomic component. Of the 18 MAD poisons, 13 consumed only part of the ergonomic component, leaving 0.03--0.45~m beyond the electrical component, and 5 were numerically below the benchmarked electrical component by 0.01--0.17~m. All 7 PRC-024-3 poisons shortened the benchmark no-trip time by 60\%. All 8 Massachusetts DER trip-threshold poisons moved the encoded threshold farther from nominal, widening the benchmark no-trip region. The remaining 2 public exact-value items were return-to-service values and are not assigned a directional consequence class here. All 45 fictional Tier-2 poisons exceeded the record limit by 20\%. 2 of these electrical-component cases (T1-025, T1-027) were among the nine $k=2$ adopted trials. These are benchmark consequence classes, not professional-engineering or physical-safety determinations.
\emph{Post-hoc coverage what-if.} Severity-prioritized D3 gave no consistent benefit: relative to hash coverage, it changed violating trials per model by -1 to +0 of 160 at 40\% and by +1 to +2 at 70\%, with Tier-3 coverage held fixed.



\section{Discussion}
Descriptively at $k=5$, retrieved poison rarely transferred its value under the conflict-abstaining prompt, but exposed Tier~1 trials usually lost the correct answer and became non-committal. These are distinct failures: non-commitment is an availability loss requiring external handling, whereas adoption is silent integrity loss. The minimal-pair construction is central: natural retrieval tends to co-retrieve the poison and its paired clean twin, Stage~B necessarily contains that twin, and locked BM25 ordering retains it whenever the poison is retrieved. Low adoption here is therefore not general protection by clean evidence. Removing permission to abstain raised adoption to 24/49 for Mistral and 8/49 for Luna. Mistral met the pre-stated material band; Luna's point estimate fell in the partial band, but its interval spans all three bands and 12/49 outputs were empty, so the Luna result is imprecise. We use these bands to describe prompt dependence, not formal causal mediation. Within the pre-specified $k=2$ condition, the post-hoc exposed-trial breakdown shows that all three original models adopted on the same nine exposed trials, from nine items; none of those reproduced contexts contained a clean same-item passage.

Stage~C narrows that association. Under the base prompt, the poison's paired clean twin yielded 0/80 adoptions and 0/80 correct answers in both locked models, while a different clean support left 15--17/80 adoptions; a lone poison yielded 80/80. Thus composition matters, but only the near-duplicate twin reliably suppressed adoption in these controls, and it did so through non-commitment rather than correct resolution. Because Stage~C did not vary the prompt, its twin result is consistent with conflict abstention but does not independently identify that instruction as the mechanism. RQ2 remains inconclusive: its fixed context necessarily contains the paired twin, so the adoption floor cannot establish general public-versus-asset resistance. Closed-book Tier~1 CRR was 0/70 for Mistral, 2/70 for Qwen, and 0/70 for Luna; Luna also returned 41/70 empty visible outputs, so its parametric prior was not adequately measured.

D3 mainly measures deterministic record coverage. Given the record key and 78/80 matching poisonable records, full coverage removes observed violations by construction; at nominal 70\%, 57/80 items were covered, so most filtering at that level is expected from coverage rather than a learned defense effect. Sparse violating items made the locked D3-vs-D1/D2 tests non-significant, so those differences remain descriptive.

Tier~3 isolates record-scope cost, whose magnitude is largely set by construction: 13/20 items (65\%) carry a generic record. Clean RAG resolved 35--39 of 40 scoped-conflict trials; the construction then illustrates how a generic SoR can flag otherwise-correct context-specific answers and why record authority must be scoped to the applicable jurisdiction or interconnection. Escalation policy and real-repository redundant evidence were not tested.

\section{Conclusion}
GridPoison-100 separates retrieval exposure, attacker adoption, non-commitment, and benchmark consequence in utility RAG. Under a conflict-abstaining prompt, one retrieved poison at $k=5$ descriptively produced mostly non-commitment rather than adoption, but this result is bounded by a minimal-pair design that often co-retrieves the paired clean twin. A locked commit-required control showed material prompt dependence for Mistral; Luna's point estimate fell in the partial band but remained imprecise. RQ2 remained inconclusive because its fixed context necessarily contained the paired twin. Stage~C showed 0/80 adoption with that twin and 80/80 with a lone poison, while a different clean support falsified the strong composition hypothesis for Mistral. At locked 70\% coverage, the keyed record check removed 5--6 of 6--8 B2 violating trials, as expected from deterministic coverage of 57/80 items; full coverage removed the remainder by construction. The deliberately generic Tier~3 records illustrate the need for jurisdiction- and interconnection-scoped authority, while generalization to real repositories remains untested.

\balance
\bibliographystyle{IEEEtran}
\bibliography{references}
\end{document}
